% ==========================================================
% Internship / Term Project Report
% BSc (Hons) Data Science and AI, IIT Guwahati
% ==========================================================

\documentclass{article}
\usepackage{spconf,amsmath,amssymb,graphicx,hyperref}
\usepackage{booktabs}
\usepackage{enumitem}

\def\x{{\mathbf x}}
\def\L{{\cal L}}

% -------------------------
% Title and Author Info
% -------------------------
\title{Next-Day Solar Irradiance Forecasting for a Tropical Island Grid:
A Leakage-Audited LSTM Encoder--Decoder Benchmarked Against Persistence}

\name{[YOUR NAME] (Roll Number: [ROLL NO])}

\address{
Term Project Report\\
Trimester [7/8/9]\\[0.5em]
B.Sc. (Honours) Data Science and Artificial Intelligence\\
Indian Institute of Technology Guwahati, India\\
\texttt{[your-email]@op.iitg.ac.in}
}

\begin{document}
\ninept
\maketitle

% ==========================================================
\begin{abstract}
% ==========================================================
Small island power systems that are adding utility-scale solar capacity need
day-ahead irradiance forecasts in order to schedule dispatchable generation, yet
they frequently lack access to a numerical weather prediction feed. I ask whether
a compact recurrent model, given nothing but strictly past-observed data, can
beat the persistence forecast that such operators fall back on. Using six years
of hourly NASA POWER reanalysis for Port of Spain, Trinidad and Tobago, I train a
two-layer LSTM encoder with a multilayer perceptron head to emit all twenty-four
hours of the next day at once, and I compare it against a twenty-four hour
persistence baseline on a held-out test year. Before recording any result I
registered the hypothesis that the model would achieve a positive forecast skill
score, with the larger gain falling in the wet season. Both parts hold. On the
2024 test year the model reduces daylight-masked RMSE from 123.35 to
90.78~W/m\textsuperscript{2}, a forecast skill of 0.264, and the seasonal split
is 0.286 wet against 0.201 dry. Because a result of this kind is worthless if the
model has seen its own answer, I supported it with twenty-seven programmatic
leakage assertions rather than with prose, and the pipeline reproduces
bit-exactly across independent runs.
\end{abstract}

% ==========================================================
\begin{keywords}
% ==========================================================
Solar irradiance forecasting, long short-term memory, time series, data leakage,
renewable energy integration
\end{keywords}

% ==========================================================
\section{Introduction}
% ==========================================================

Trinidad and Tobago has committed to sourcing thirty percent of its electricity
from renewables by 2030 under its Paris Agreement contribution
\cite{gott2021tt}, and its first utility-scale solar facility is now feeding a
grid that has run on natural gas for its entire history. That transition creates
an operational problem that did not previously exist. A gas turbine produces what
it is instructed to produce; a solar farm produces what the sky permits. To keep
supply matched to demand, a system operator must decide the evening before how
much dispatchable generation to commit for the following day, and the quality of
that decision depends directly on how well tomorrow's irradiance can be
anticipated.

The forecasting literature is mature and largely agrees on what works
\cite{yang2018history, inman2013solar}, but almost all of the strongest methods
assume access to a numerical weather prediction (NWP) feed. That assumption is
comfortable for a European transmission operator and much less so for a small
island utility, for which an NWP subscription is a procurement decision with a
price attached. The operator without one falls back on persistence: assume
tomorrow looks like today. In the tropics this is a weaker crutch than it sounds,
because convective cloud makes consecutive days only loosely related.

My interest is therefore in the gap between those two positions. I ask what a
learned model can extract from observational history alone, with no forecast
information of any kind, and whether that is enough to justify displacing
persistence. This framing has a consequence that shapes the entire design: the
model may see tomorrow's calendar, because a clock requires no forecast, but it
may not see tomorrow's weather, because that would presuppose the very feed the
problem assumes is absent.

Section~2 states the problem and my objectives. Section~3 details the
methodology, including the architecture and the leakage audit. Section~4 presents
the experimental results, Section~5 discusses and interprets them, and Section~6
concludes with future work.

% ==========================================================
\section{Problem Statement and Objectives}
% ==========================================================

\subsection{Problem Statement}

Let $y_t$ denote global horizontal irradiance (GHI) in W/m\textsuperscript{2} at
hour $t$, and let $\mathbf{w}_t$ denote a vector of concurrently observed weather
variables. A forecast is issued at local midnight on each day, which I take as
the forecast origin $t$. The task is to produce
$\hat{y}_{t+1}, \dots, \hat{y}_{t+24}$ using only information available at time
$t$.

The binding constraint is what "available" means. The model receives the previous
twenty-four hours of observed irradiance and weather, $\{y_{t-23:t},
\mathbf{w}_{t-23:t}\}$, together with deterministic calendar features for the
target hours. It does not receive any observation or forecast of the weather
during $t+1$ to $t+24$. This is deliberately harsher than the closest prior work:
Qing and Niu \cite{qing2018hourly} forecast day-ahead hourly irradiance for an
island site using an LSTM but supply the target day's weather forecasts as input.
Withholding them is the specific question I am asking.

\subsection{Objectives}

\begin{enumerate}[noitemsep,topsep=2pt]
    \item To assemble and validate a six-year hourly meteorological record for
      Port of Spain, and to verify it is physically plausible before modelling.
    \item To implement a twenty-four hour persistence forecast as the reference
      against which skill is measured.
    \item To design and train an LSTM encoder--decoder that forecasts all
      twenty-four hours of the next day using only past-observed inputs.
    \item To prove, by programmatic assertion rather than by inspection, that no
      target information reaches the model.
    \item To evaluate on a held-out year using daylight-masked error metrics and
      the forecast skill score, disaggregated by season and hour of day.
    \item To quantify what a perfect weather forecast would be worth, by training
      an upper-bound variant that is given the target day's observed weather.
\end{enumerate}

Before running any model I recorded the hypothesis that the LSTM would achieve a
skill score above zero on the held-out test year, with the largest gains in
wet-season months. Registering it in advance is what makes the seasonal result in
Section~4 a confirmation rather than an observation found after the fact.

% ==========================================================
\section{Methodology}
% ==========================================================

\subsection{Data and Features}

I used the NASA POWER hourly reanalysis \cite{sparks2018nasapower} at
$10.65^\circ$N, $61.50^\circ$W for 2019--2024 inclusive, in local solar time:
52{,}608 hours with six variables, namely all-sky surface shortwave downward flux
(taken as GHI), air temperature, relative humidity, wind speed, cloud amount and
precipitable water. POWER serves hourly GHI in Wh/m\textsuperscript{2}, which is
numerically equal to the mean W/m\textsuperscript{2} over that hour, so no
conversion factor is applied anywhere.

To each hour I added five calendar features. Hour of day and day of year are
cyclical, so encoding them as integers would place 23:00 and 00:00 at opposite
ends of a range they are in fact adjacent to. I encoded each as a sine--cosine
pair,
\begin{equation}
    \left(\sin\tfrac{2\pi h}{24},\ \cos\tfrac{2\pi h}{24},\
          \sin\tfrac{2\pi d}{365.25},\ \cos\tfrac{2\pi d}{365.25}\right),
\end{equation}
which preserves adjacency, plus a binary wet-season indicator for June through
December. These five are the only features the model is permitted to know for
tomorrow.

\subsection{Windowing, Splits and Scaling}

Each training example pairs an encoder input of shape $(24, 11)$ with a decoder
input of shape $(24, 5)$ and a target of shape $(24,)$. Splits are chronological
by calendar year, never random: 2019--2022 train, 2023 validation, 2024 test. A
window joins a split only if its encoder input \emph{and} its target both fall
entirely inside that split's years, which costs a handful of windows at each year
boundary but guarantees the splits share no hour. Training uses stride~1, giving
35{,}017 windows; validation and test use stride~24 at midnight origins, giving
363 and 364 windows, so that evaluation measures the operational task of one
forecast issued per day. Standardisation statistics are fitted on training rows
only, and the fitted index range is persisted so that the claim can be audited
rather than assumed.

\subsection{Model Architecture}

The encoder is a two-layer LSTM \cite{hochreiter1997long} with 64 hidden units,
dropout 0.2 between layers \cite{srivastava2014dropout}. At each step it computes
\begin{align}
    \mathbf{i}_t &= \sigma(W_i \x_t + U_i \mathbf{h}_{t-1} + b_i), \\
    \mathbf{f}_t &= \sigma(W_f \x_t + U_f \mathbf{h}_{t-1} + b_f), \\
    \mathbf{o}_t &= \sigma(W_o \x_t + U_o \mathbf{h}_{t-1} + b_o), \\
    \mathbf{c}_t &= \mathbf{f}_t \odot \mathbf{c}_{t-1}
                    + \mathbf{i}_t \odot \tanh(W_c \x_t + U_c \mathbf{h}_{t-1} + b_c), \\
    \mathbf{h}_t &= \mathbf{o}_t \odot \tanh(\mathbf{c}_t),
\end{align}
where $\odot$ is the elementwise product. The gating is what makes the recurrence
appropriate here: the forget gate lets the network retain a persistent cloud
regime across the full day while discarding hour-to-hour fluctuation.

I take the final hidden state $\mathbf{h}_{24}$ of the last layer as a summary of
what the sky has been doing, concatenate it with the flattened calendar block
$\mathbf{D} \in \mathbb{R}^{24\times5}$ describing which twenty-four hours are
being asked about, and pass the result through a single hidden layer of 128 ReLU
units:
\begin{equation}
    \hat{\mathbf{y}} = W_2\,\mathrm{ReLU}\!\left(W_1
      \begin{bmatrix}\mathbf{h}_{24} \\ \mathrm{vec}(\mathbf{D})\end{bmatrix}
      + b_1\right) + b_2, \quad \hat{\mathbf{y}} \in \mathbb{R}^{24}.
\end{equation}
Concatenation rather than conditioning keeps the two questions separable, so the
calendar path can learn the clear-sky envelope while the recurrent path learns
the departure from it. All twenty-four hours are emitted simultaneously. A
recursive decoder \cite{sutskever2014sequence} would feed its own hour-one
prediction into hour two and compound error along the horizon, and it would
require a GHI input for tomorrow that by construction does not exist. The model
holds 79{,}768 trainable parameters.

\subsection{Training}

I trained with Adam \cite{kingma2015adam} at learning rate $10^{-3}$, batch size
64, plain MSE over all twenty-four hours, and early stopping on validation loss
with patience 10. Training halted at epoch 15 and restored epoch 5. Execution is
CPU-only with deterministic algorithms enabled, because cuDNN's LSTM kernels are
not deterministic; consequently two independent clean runs reproduce all 312
metric rows bit-exactly, not merely to three decimal places.

\subsection{Validation Gates}

A forecasting result is only as trustworthy as the evidence that the model never
saw its answer, and prose assurances are not evidence. I therefore built three
gates that halt the pipeline on failure. Gate~0 checks the raw download for
physical plausibility before anything else runs: row count, night mean below
20~W/m\textsuperscript{2}, midday mean within $[300, 1100]$, maximum below 1200,
and a per-column missing census. Gate~1 runs the entire preprocess, windowing,
training and evaluation path on a synthetic clipped sinusoid whose answer is
known by construction, so that a failure can be attributed to the pipeline rather
than to the model. Gate~2 is the leakage audit: twenty-seven assertions verifying
not only that every encoder timestamp precedes every target timestamp, but that
the encoder tensor \emph{equals} the scaled observations at those timestamps and
the target tensor equals the scaled GHI at its own, that the decoder carries the
five calendar features and nothing else, that the three splits share no hour, and
that the scaler's recorded fit index terminates before the validation year opens.
Checking values rather than shapes is what makes this a proof rather than a
gesture. All three gates pass.

\subsection{Baseline and Metrics}

The reference forecast is twenty-four hour persistence,
$\hat{y}_{t+k} = y_{t+k-24}$, which for a midnight origin means yesterday
repeated. I report MAE, RMSE and MAPE, and the forecast skill score
\cite{marquez2013proposed}
\begin{equation}
    s = 1 - \frac{\mathrm{RMSE}_{\text{model}}}{\mathrm{RMSE}_{\text{ref}}},
\end{equation}
where $s = 0$ means no improvement on persistence and $s = 1$ means perfect
prediction. All headline metrics are computed only on hours where observed GHI
exceeds 20~W/m\textsuperscript{2}. Roughly half of any twenty-four hour window is
night, when GHI is zero and every model is correct for free; including those
hours would inflate both models' apparent accuracy and compress the difference
between them. Predictions are clipped at zero before scoring, since irradiance
cannot be negative and persistence, being composed of real observations, can
never violate that floor.

% ==========================================================
\section{Experiments and Results}
% ==========================================================

\subsection{Experimental Setup}

All results below are computed on 2024, a year touched exactly once, after all
design decisions were fixed against the 2023 validation year. The test year
contributes 364 daily forecasts, of which 4{,}229 of 8{,}736 hours survive the
daylight mask. The pipeline runs end to end in roughly 75 seconds on a laptop CPU
in Python 3.12 with PyTorch, scikit-learn, pandas and Matplotlib.

\begin{table}[t]
\centering
\caption{Test-year (2024) performance, daylight-masked. Skill is measured
against persistence. The upper bound is not a forecast; it is given the target
day's observed weather.}
\label{tab:main}
\footnotesize
\begin{tabular}{lrrrr}
\toprule
Model & MAE & RMSE & MAPE & Skill \\
      & (W/m\textsuperscript{2}) & (W/m\textsuperscript{2}) & (\%) & \\
\midrule
Persistence        & 77.85 & 123.35 & 21.64 & 0.000 \\
\textbf{LSTM}      & \textbf{62.82} & \textbf{90.78} & \textbf{21.01} & \textbf{0.264} \\
\midrule
\emph{Upper bound} & \emph{39.34} & \emph{62.11} & \emph{15.17} & \emph{0.496} \\
\bottomrule
\end{tabular}
\end{table}

\subsection{Results}

Table~\ref{tab:main} reports the headline comparison. The LSTM reduces
daylight-masked RMSE by 26.4 percent relative to persistence, corresponding to a
forecast skill of 0.264, and reduces MAE by 19.3 percent. Both parts of the
registered hypothesis are confirmed: skill is positive, and the seasonal
breakdown in Fig.~\ref{fig:seasonal} gives 0.286 in the wet season against 0.201
in the dry.

\begin{figure}[t]
    \centering
    \includegraphics[width=\linewidth]{figures/fig4_seasonal.png}
    \caption{Daylight-masked RMSE by month of the test year. The wet season
    (June--December) is shaded. Persistence degrades sharply where the LSTM does
    not, which is where the skill is earned.}
    \label{fig:seasonal}
\end{figure}

\begin{figure}[t]
    \centering
    \includegraphics[width=\linewidth]{figures/fig2_cloudy_day.png}
    \caption{The most heavily clouded day of the test year (2024-08-24).
    Persistence replays a brighter previous day and overshoots badly; the LSTM,
    having read the preceding twenty-four hours, forecasts a suppressed profile.}
    \label{fig:cloudy}
\end{figure}

Fig.~\ref{fig:seasonal} shows where the aggregate number comes from. Persistence
error peaks in July at roughly 177~W/m\textsuperscript{2} while the LSTM stays
near 118, and the two models are nearly indistinguishable in February and March.
Figs.~\ref{fig:cloudy} and~\ref{fig:clear} show the same contrast at the scale of
a single day, and are best read as a pair: on the clearest day of the test year
persistence is a perfectly reasonable forecast, and on the most clouded day it is
not. Fig.~\ref{fig:hour} disaggregates by hour and shows that error for both
models is concentrated around solar noon, when the absolute magnitude of
irradiance, and therefore the scope to be wrong about it, is greatest.
Fig.~\ref{fig:loss} records the training history: validation loss bottomed at
epoch 5 and rose steadily thereafter while training loss kept falling, so early
stopping restored epoch 5 and discarded the ten epochs of overfitting that
followed.

\begin{figure}[t]
    \centering
    \includegraphics[width=\linewidth]{figures/fig1_clear_day.png}
    \caption{The clearest day of the test year (2024-09-05). Under a stable sky
    both models track the observed profile closely, which is why the aggregate
    skill in Table~\ref{tab:main} is not driven by days like this one.}
    \label{fig:clear}
\end{figure}

%% fig:loss is defined after fig:hour so that float numbering follows the order
%% in which the figures are referred to in the text.

\begin{figure}[t]
    \centering
    \includegraphics[width=\linewidth]{figures/fig3_error_by_hour.png}
    \caption{Daylight-masked MAE by hour of day. Both models are worst near solar
    noon; the LSTM's advantage is widest there and narrows towards the terminators.}
    \label{fig:hour}
\end{figure}

\begin{figure}[t]
    \centering
    \includegraphics[width=\linewidth]{figures/fig6_loss_curve.png}
    \caption{Training and validation loss. The divergence after epoch 5 is
    textbook overfitting; early stopping fired at epoch 15 and the epoch-5
    weights were restored.}
    \label{fig:loss}
\end{figure}

% ==========================================================
\section{Discussion}
% ==========================================================

\textbf{Where the skill comes from.} Persistence carries yesterday's cloud field
forward unchanged, which is an excellent assumption under a stable synoptic
regime and a poor one under convection. The wet-season result is therefore not a
quirk of my model but a statement about the baseline: persistence degrades from
94.51 to 140.15~W/m\textsuperscript{2} between dry and wet seasons, a 48 percent
deterioration, while the LSTM degrades from 75.55 to 100.13, a 33 percent
deterioration. The learned model earns its keep precisely in the regime where the
naive assumption breaks, which is the regime an operator most needs help with.
The mechanism is visible in Fig.~\ref{fig:cloudy}: given an ambiguous encoder
history the model hedges towards the seasonal clear-sky envelope it has learned
from the calendar features, whereas persistence has no capacity to hedge at all.

\textbf{A metric that disagrees.} MAPE separates the two models by almost
nothing, 21.64 against 21.01 percent, while RMSE separates them by 26 percent.
This is not a contradiction but a property of MAPE: it divides by the observed
value, so hours barely above the 20~W/m\textsuperscript{2} daylight threshold
dominate the average, and those hours near dawn and dusk are exactly where the
two models most nearly agree. I report MAPE because it is conventional, but I
would not use it to select a model for this task, and I lead on RMSE and skill
score for that reason. An unmasked MAPE is undefined here in any case, because
night GHI is exactly zero.

\textbf{What a weather feed would be worth.} The upper-bound variant, handed the
target day's observed weather but never its observed GHI, reaches a skill of
0.496 against the primary model's 0.264. It is not a forecast and I never present
it as one. Its value is that it converts an abstract question into a number: a
hypothetical perfect NWP feed would roughly double the gain over persistence, and
that is the figure to weigh against the cost of a subscription. It also bounds my
own result usefully, showing that roughly half of the achievable improvement is
available with no forecast information whatsoever.

\textbf{Limitations.} Four are worth stating plainly. First, this is a single
site, and skill scores do not transfer across climates, so nothing here should be
read as a claim about tropical sites in general. Second, NASA POWER is
satellite-derived reanalysis rather than a ground pyranometer; it is internally
consistent and appropriate for a proof of concept, but a deployed system would
need validation against station data. Third, I compared against persistence
alone. A clear-sky-index formulation \cite{yang2020choice} or a gradient-boosted
tree on the same features would be a more demanding reference, and I would expect
them to close part of the gap. Fourth, the model is deterministic and emits a
point forecast, whereas an operator sizing a reserve margin needs an interval.

% ==========================================================
\section{Conclusion and Future Work}
% ==========================================================

I set out to test whether a compact LSTM restricted to strictly past-observed
data could beat persistence at next-day GHI forecasting for Port of Spain. It
can: forecast skill of 0.264 on a held-out year, with the pre-registered
wet-season advantage confirmed at 0.286 against 0.201. The result rests on
twenty-seven programmatic leakage assertions and a bit-exactly reproducible
pipeline, which I regard as the more durable contribution, since a skill score is
only as credible as the audit behind it.

Three extensions follow directly. The most valuable is a stronger reference:
adding a clear-sky-index persistence model and a gradient-boosted baseline would
establish how much of the 0.264 is attributable to recurrence rather than to
merely having any learned model. The second is probabilistic output, via quantile
regression or a pinball loss, which turns the forecast into the reserve-sizing
input an operator actually needs. The third is spatial extension to several sites
across the two islands, which would test transferability and open the possibility
of a single model conditioned on location.

% ==========================================================
\section{Artifacts and Demonstrations}
% ==========================================================

\begin{itemize}[noitemsep,topsep=2pt]
    \item \textbf{Code repository:} \url{[GITHUB URL]} --- full pipeline,
      reproducible with a single command. Contains \texttt{DECISIONS.md}, a
      numbered record of every discretionary choice with its rationale, and
      \texttt{results/validation\_report.txt}, the full output of all three gates.
    \item \textbf{Video presentation:} \url{[YOUTUBE URL]}
\end{itemize}

AI-assisted tools were used for structuring and formatting this report. The
research design, implementation, analysis and conclusions are my own.

% ==========================================================
% References
% ==========================================================
\vfill\pagebreak
{\fontsize{8}{9}\selectfont
\bibliographystyle{IEEEbib}
\bibliography{refs}
}

\end{document}
